\chapter{测试功能与算法设计}

\section{测试功能总体流程}

通信与基础控制模块已经能够发送控制帧、读取反馈样本并判断连接状态。测试功能设计接着要处理的问题，是如何把这些反馈数据转化为可评价的性能指标。伺服系统测试平台研究通常按功能测试、性能测试和测试流程展开方案设计，工业测试系统研究也把测试流程、数据管理和可靠性验证放在同一条链路中讨论\cite{li2014servotest,xu2021industry4test,liu2022industrialsoftware}。结合舵机测试场景，软件将测试功能分为通信性能测试、控制性能预测试、时域性能测试和频域性能测试四类，前端参数面板也按这一顺序组织。

四类测试按执行顺序排布。通信性能测试先确认串口链路、协议解析和反馈接收能否支撑后续任务；控制性能预测试随后估计噪声标准差、保守调节时间和平均运动速度；时域性能测试利用这些参数确定采样窗口、重复次数和测试路径，并从位置响应过程中提取动态指标；频域性能测试则按照噪声水平和调节时间确定每个频率点的采样规模，再通过相关法估计幅值和相位。基础通信能力、统计参数估计、时域响应分析和频域响应分析由此进入同一套执行流程。

桌面测控软件通常要在控制、采集、显示和结果处理之间建立稳定的数据边界\cite{zhang2017qtflowmeter,li2014servotest}。本系统中的四类测试核心都通过统一运行时接口获取控制发送、反馈读取、延时和当前时间等能力，并通过进度、统计、完成和错误事件向界面层回传过程状态。测试算法据此聚焦于数据采集和指标计算，界面层保留参数输入、进度显示和结果展示职责。围绕单对象控制实验平台的研究也表明，控制执行、反馈采集和上位机显示的一体化组织，有助于降低实验过程中的状态割裂和操作跳转\cite{lin2025dtcplatform,zhou2020pmsmexperiment,zhang2016bldcexperiment}。

\section{通信性能测试设计}

通信性能测试用于评价当前协议和波特率配置下的串口链路。它关注控制帧能否发出、反馈帧能否收到，以及反馈间隔是否稳定。发送异常、反馈丢失和反馈时间间隔波动都会影响后续动态指标，因此系统把通信性能测试放在四类测试的起点。针对 RS422 链路的误码测试研究表明，发送、接收与错误统计需要作为统一流程组织，而分包和缓存机制会直接影响连续测试中的稳定性\cite{hu2026rs422ber,wu2025serialpacket}。

测试过程包含发送测试、接收测试和联合测试三种模式。发送测试统计控制帧发送成功次数，接收测试统计有效反馈帧接收次数并记录反馈时间间隔，联合测试在同一任务内同时统计发送和接收结果。设发送成功次数为 $n_s$、发送总次数为 $N_s$，接收成功次数为 $n_r$、接收总次数为 $N_r$，则发送成功率和接收成功率分别为
\begin{equation}
  P_s=\frac{n_s}{N_s}\times 100\%,\qquad
  P_r=\frac{n_r}{N_r}\times 100\%.
\end{equation}

反馈时间间隔用于描述链路在连续测试中的时间稳定性。系统在接收过程中在线更新均值和方差，不等到任务结束后再处理全量间隔样本。设第 $k$ 个反馈时间间隔为 $x_k$，当前均值为 $\mu_k$，平方差累计量为 $M_{2,k}$，则更新关系为
\begin{equation}
  \mu_k=\mu_{k-1}+\frac{x_k-\mu_{k-1}}{k},
\end{equation}
\begin{equation}
  M_{2,k}=M_{2,k-1}+(x_k-\mu_{k-1})(x_k-\mu_k).
\end{equation}

由此得到时间间隔标准差和 $95\%$ 置信区间：
\begin{equation}
  \sigma_k=\sqrt{\frac{M_{2,k}}{k}},\qquad
  CI_{95}=\left[\mu_k-1.96\frac{\sigma_k}{\sqrt{k}},\mu_k+1.96\frac{\sigma_k}{\sqrt{k}}\right].
\end{equation}

在线统计使通信测试可以随着反馈到达持续更新指标，并将均值、标准差、最小值、最大值和置信区间同步到界面。测试过程中，每个任务都绑定唯一任务标识，进度事件记录当前阶段、已发送次数和已接收次数，错误事件区分测试失败和任务取消。通信质量先在这一环节确认，可以减少后续算法把链路异常误判为舵机响应异常的风险。

\begin{figure}[H]
  \centering
  \includegraphics[width=0.68\textwidth]{ch5-comm-test.png}
  \caption{通信性能测试结果界面}
  \label{fig:comm-test-result}
\end{figure}
\FloatBarrier

\section{控制性能预测试设计}

控制性能预测试承担参数过渡作用。时域测试和频域测试都要用到噪声水平、调节时间和响应速度，完全依赖经验填写会降低结果解释的稳定性。实验设计研究强调，测试方案需要在参数可信度与执行成本之间取得平衡\cite{optimisation2020intro}。系统在正式性能测试前先做预测试，把静止噪声采集和阶跃响应测量作为后续参数来源。

预测试由静止噪声采集和阶跃响应测试两部分组成。静止噪声采集阶段，系统让舵机保持在固定目标位置，并在设定采样时间内读取反馈样本。设反馈样本数为 $M$，样本标准差为 $s$，置信水平参数为 $\alpha$，则噪声标准差采用保守估计：
\begin{equation}
  \sigma_n=s\sqrt{\frac{M-1}{\chi^2_{\alpha,M-1}}}.
\end{equation}
其中，$\chi^2_{\alpha,M-1}$ 为自由度 $M-1$ 的卡方分布分位数。该估计把有限样本带来的不确定性计入噪声参数，后续时域和频域测试在计算采样规模时使用更保守的噪声水平。

阶跃响应阶段用于估计调节时间和平均运动速度。系统采用连续阶跃路径
\begin{equation}
  25^\circ \rightarrow 0^\circ \rightarrow 25^\circ \rightarrow 0^\circ \rightarrow 25^\circ \rightarrow 0^\circ
\end{equation}
对舵机进行多次正反方向切换。对于第 $i$ 次阶跃，设起始位置和目标位置分别为 $\theta_{i-1}$ 和 $\theta_i$，调节时间为 $T_{s,i}$，则该次阶跃对应的平均速度为
\begin{equation}
  v_i=\frac{|\theta_i-\theta_{i-1}|}{T_{s,i}}.
\end{equation}

预测试输出的保守调节时间和平均速度定义为
\begin{equation}
  T_s=\max_i T_{s,i},\qquad
  \bar{v}=\frac{1}{L}\sum_{i=1}^{L}v_i,
\end{equation}
其中，$L$ 为有效阶跃次数。保守调节时间用于后续测试中的等待时间和单点最大采样时间，平均速度用于估计测试路径执行耗时。阶段性实机验证表明，预测试流程能够给出后续测试所需参数，但该结果只说明当前工况下可用，不能直接推广到所有场景。

\begin{figure}[H]
  \centering
  \includegraphics[width=0.94\textwidth]{ch5-pretest.png}
  \caption{控制性能预测试结果界面}
  \label{fig:pretest-result}
\end{figure}
\FloatBarrier

\section{时域性能测试设计}

时域性能测试直接面向位置响应过程，目标是在有限测试时间内获得位置误差、动态响应和不同运动方向下的一致性信息。该测试从预测试读取噪声标准差 $\sigma_n$、保守调节时间 $T_s$ 和平均速度 $\bar{v}$，再由操作人员设置置信水平、容许误差和离散点数。参数来源由预测试提供后，时域测试中等待时间和重复次数的设置不再完全依赖人工经验。

测试位置在舵机物理范围 $[-25^\circ,25^\circ]$ 内均匀分布。设离散点数为 $M$，则第 $j$ 个测试位置为
\begin{equation}
  p_j=-25+\frac{50(j-1)}{M-1},\qquad j=1,2,\ldots,M.
\end{equation}

容许误差 $\delta$ 和噪声标准差 $\sigma_n$ 决定统计窗口长度。设 $z_{1-\alpha/2}$ 为标准正态分位数，则窗口长度为
\begin{equation}
  N=\left\lceil \left(\frac{z_{1-\alpha/2}\sigma_n}{\delta}\right)^2\right\rceil.
\end{equation}

重复次数 $R$ 通过 $t$ 分布临界值迭代获得，并设置上限以避免极端参数把测试时间拉得过长。时域测试的采样规模由噪声水平、置信要求和容许误差共同决定，固定次数只作为边界约束。

测试路径需要兼顾覆盖性和耗时。全组合遍历会使总点数随 $M$ 快速增加，因此系统采用两阶段路径：第一阶段完成全量程正向扫描和反向扫描，保证每个相邻位置和两个运动方向都被覆盖；第二阶段只在 $K=\min(5,M)$ 个代表位置进行重复，以获得重复一致性数据。总测试点数为
\begin{equation}
  L=2(M-1)+(R-1)2K.
\end{equation}
该路径保留了位置覆盖、方向比较和重复一致性分析所需样本，同时控制总测试时间。路径组织以覆盖需求与测试耗时之间的平衡为目标，避免全组合遍历导致测试时间过长。

单个测试点执行时，系统发送目标位置控制量，并在 $T_s+3000$ ms 的最大时长内采集反馈曲线。反馈曲线以时间和角度构成样本序列，系统使用滑动均值判断是否进入目标位置的容差带。若最近若干样本均值满足
\begin{equation}
  |\bar{y}(t)-p_j|\le \delta
\end{equation}
且采样时间已经超过预测试得到的保守调节时间，则该测试点可以结束。响应曲线随后用于提取调节时间、超调量、上升时间、峰值时间、稳态误差、延迟时间和振荡次数等指标。上升时间由 $10\%$ 到 $90\%$ 阈值穿越时间给出，稳态误差由末段反馈均值与目标位置之差给出，振荡次数由相对稳态偏差的过零次数折算得到。

时域测试结果按重复一致性、位置均匀性、正反向偏差和步长相关性四种方式分组，分别观察调节时间、稳态误差、回程差异和不同阶跃长度下的超调情况。动态指标仍从有效测试点中提取上升时间、峰值时间、超调百分比、调节时间、稳态误差、延迟时间和振荡次数。时域测试既给出单点响应过程，也给出跨位置、跨方向和跨步长的整体评价。

\begin{figure}[H]
  \centering
  \includegraphics[width=0.68\textwidth]{ch5-time-domain.png}
  \caption{时域性能测试结果界面}
  \label{fig:time-domain-result}
\end{figure}
\FloatBarrier

\section{频域性能测试设计}

频域性能测试从频率响应角度评价舵机对正弦激励的跟随能力。与时域测试相同，频域测试也读取预测试得到的噪声标准差和保守调节时间；不同之处在于，采样规模要按扫频范围和估计精度来定，再从反馈信号中估计输出幅值和相位。该测试结果主要以幅频曲线、相频曲线和若干频域指标呈现。

扫频点按对数间隔生成。设起始频率为 $f_{\min}$，终止频率为 $f_{\max}$，频率点数为 $K$，则第 $k$ 个扫频点为
\begin{equation}
  f_k=f_{\min}\left(\frac{f_{\max}}{f_{\min}}\right)^{\frac{k-1}{K-1}},\qquad k=1,2,\ldots,K.
\end{equation}
对数扫频能够在低频段保留足够分辨率，也避免高频段点数过少。

采样点数由幅值估计精度、相位估计精度和频率覆盖共同决定。设输入正弦幅值为 $A$，噪声标准差为 $\sigma_n$，幅值容许误差为 $\delta_A$，相位容许误差为 $\delta_\phi$，则幅值估计约束对应的样本数为
\begin{equation}
  N_A=\left\lceil \left(\frac{z\sigma_n\sqrt{2}}{\delta_A}\right)^2\right\rceil,
\end{equation}
相位估计约束对应的样本数为
\begin{equation}
  N_{\phi}=\left\lceil \left(\frac{z\sigma_n}{A\delta_{\phi}}\right)^2\right\rceil.
\end{equation}
为保证相关法具有足够周期覆盖，系统还要求每个频率点至少覆盖 3 个完整周期。若控制发送间隔为 $T_{\mathrm{cmd}}$，频率点周期为 $T_k=1/f_k$，则实际采样点数取
\begin{equation}
  N_k=\max\left(N_A,N_{\phi},\left\lceil\frac{3T_k}{T_{\mathrm{cmd}}}\right\rceil,4\right).
\end{equation}

完成采样后，系统采用相关法估计输出幅值和相位。设反馈样本为 $y_k$，对应时间为 $t_k$，角频率为 $\omega$，则同相分量和正交分量为
\begin{equation}
  I=\frac{2}{N}\sum_{k=1}^{N}y_k\sin(\omega t_k),\qquad
  Q=\frac{2}{N}\sum_{k=1}^{N}y_k\cos(\omega t_k).
\end{equation}
输出幅值 $B$、增益 $G$、分贝增益 $G_{dB}$ 和相位 $\phi$ 分别为
\begin{equation}
  B=\sqrt{I^2+Q^2},\qquad
  G=\frac{B}{A},\qquad
  G_{dB}=20\log_{10}G,\qquad
  \phi=\operatorname{atan2}(Q,I).
\end{equation}

相关法利用已知激励频率进行正交解调，比直接读取峰值更适合含噪反馈序列。Cajetinac 等对执行机构频率特性辨识的研究也说明，频域评估需要同时处理激励、采样和计算之间的耦合关系\cite{cajetinac2012pwm}。系统得到每个频率点的增益和相位后，再对曲线进行 3 点移动平均平滑，并从平滑曲线上提取 $-3$ dB 带宽、谐振频率、谐振峰值、相位裕度和增益裕度。$-3$ dB 带宽由增益曲线首次降至 $-3$ dB 的位置线性插值得到；谐振频率为平滑增益峰值对应频率；相位裕度和增益裕度分别依据增益穿越和相位穿越位置估计。频域测试最终形成伯德图和指标摘要，为舵机动态响应能力提供与时域测试互补的评价。

\begin{figure}[H]
  \centering
  \includegraphics[width=0.94\textwidth]{ch5-frequency-domain.png}
  \caption{频域性能测试结果界面}
  \label{fig:frequency-domain-result}
\end{figure}
\FloatBarrier

\section{测试结果组织与参数传递}

四类测试的评价目标不同，但进入程序后的数据形态保持一致。每项测试都从参数面板发起，测试核心通过统一运行时接口发送控制量、读取最近反馈样本，并在执行过程中回传进度和统计事件。任务完成后，结构化结果写入共享结果状态；数据显示区域据此切换通信成功率、预测试参数、时域动态指标或频域极值矩阵，数据保存区域则从同一状态中汇总可保存内容。

参数传递按测试顺序推进。通信测试先给出链路状态，预测试再得到噪声标准差、保守调节时间和平均速度；时域测试用这些量计算窗口长度、重复次数和执行路径，频域测试据此确定采样点数和等待时间。结果保存后，历史数据会重新写回共享结果状态，图表和指标区域可以复现已有测试记录。一次测试因此会留下可复查、可对比的数据对象，超出当次界面显示的临时状态。

参数传递还要说明适用边界。预测试输出的 $\sigma_n$ 表示静止反馈段的角度波动，主要影响统计窗口和频域幅相估计所需样本量；$T_s$ 表示已测阶跃中的保守收敛时间，主要影响单点等待时间和扫频稳定时间；$\bar{v}$ 表示阶跃过程的平均运动速度，用于估计时域路径执行耗时。三类参数都服务于测试规模估计，不能脱离舵机物理角度范围、控制发送周期和测试时长上限单独使用。

当前实现中，时域和频域测试面板会先读取最近一次预测试结果。若没有可用结果，界面使用默认值，操作人员也可以直接修改。后端测试核心在任务启动前检查置信水平、容许误差、噪声标准差、调节时间、平均速度、扫频范围和激励幅值等参数。非有限值、非正噪声、负调节时间、非正平均速度，以及超过舵机物理范围的容许误差或噪声估计，都会被判定为非法配置并通过错误事件反馈到界面。控制量进入协议帧之前还会按对应协议范围截断，避免越界值直接写入串口帧。

采样规模同样设置了工程边界。时域测试在计算重复次数时设置下限和上限，防止极小噪声或过高精度要求使测试时间失去约束；频域测试在幅值误差、相位误差和周期覆盖之间取较大采样量，并设置最少样本数。这些规则界定了当前软件对参数传递的保护范围。高噪声环境、慢响应舵机和长时间连续测试下的参数敏感性仍需要后续专门实验补充。

\section{本章小结}

本章在通信协议与基础控制模块基础上，说明四类测试功能如何进入软件执行过程。通信性能测试使用成功率和反馈时间间隔判断链路质量；控制性能预测试通过噪声估计和阶跃响应测量给正式测试提供统计参数；时域性能测试根据预测试参数生成位置路径并提取响应指标；频域性能测试通过对数扫频和相关法估计增益、相位及伯德图指标。

这四类测试沿同一测试链路递进展开。通信测试提供链路前提，预测试提供参数依据，时域和频域测试分别从时间响应和频率响应角度评价舵机性能。参数传递部分进一步说明默认值、人工修正、后端校验和采样规模约束之间的关系，结果组织和保存机制则保证测试输出能够被查看、复查和留档。第六章的系统实现与测试验证将以这些功能边界和指标输出为依据。
